\chapter{LITERATURE REVIEW}
\hspace{0.5cm}
The evaluation of software quality metrics and vulnerability localization models remains a core research frontier within software engineering and program security ecosystems. Historically, code analysis strategies have been split into two primary areas: rule-based static application security testing (SAST) and learning-based deep representation networks. While each methodology provides standalone analytical value, both exhibit significant design limitations when applied to the unique structural constraints and execution paradigms of the Rust programming language.

Traditional static analysis toolsets, exemplified by industry standard engines such as Fortify and SonarQube, rely heavily on deterministic pattern matching, text regex scanning, and pre-formulated rule databases to capture insecure code constructs. In the context of the Rust programming ecosystem, this paradigm is executed by the official compiler linter, Clippy, which relies on predefined abstract syntax checks to flag unidiomatic expressions. While highly effective at identifying shallow formatting irregularities or basic syntactic variations, these tools fall short when assessing deep logical vulnerabilities or cross-layer behavioral bottlenecks. Because linear token-matching mechanisms lack comprehensive data-flow tracking, they often produce high false-positive rates, miss complex structural defects that span conditional execution pathways, and fail to provide developers with actionable architectural reasoning.

To overcome the challenges of syntax-dependent code checking, compiler research has shifted toward structural graph abstractions that capture program behavior independently of source text. The foundational breakthrough in this domain was established by Yamaguchi et al. (2014), who formalized the concept of the Code Property Graph (CPG). By merging three distinct compiler representations—the grammatical Abstract Syntax Tree (AST), the execution-path-bound Control Flow Graph (CFG), and the data-dependence-tracking Program Dependence Graph (PDG)—into a single queryable graph database framework, the CPG transformed static analysis into a graph traversal problem. Tools built on this paradigm, such as the open-source analyzer \textbf{Joern}, allow security engineers to trace complex data propagation flows from input sources to vulnerable operation sinks. However, these traditional graph-based approaches remain strictly structural and deterministic. Their accuracy depends entirely on manual query development using specialized domain languages, leaving them unable to adapt to novel, context-dependent architectural flaws. Furthermore, classical CPG models are optimized for C/C++ memory models and lack the specialized semantic nodes needed to reason about Rust’s unique ownership, move, and reference borrowing rules.

Concurrently, the integration of deep learning and Natural Language Processing (NLP) has led to the development of massive transformer architectures trained on vast code corpora. Pre-trained language models for code, such as CodeBERT and GraphCodeBERT, have demonstrated strong capabilities across diverse software tasks, including vulnerability classification, defect prediction, and code summarization. Despite these advances, pure deep learning models face severe limitations when operating directly on raw high-level source code. As highlighted by Chakraborty et al. (2022) in their evaluation "\textit{Deep Learning Based Vulnerability Detection: Are We There Yet?}", token-based language models often lack true compiler-level comprehension. They are highly sensitive to superficial syntax modifications, variable renaming, and coding styles, which leads to unpredictable accuracy drops and high false-negative rates when processing subtle logic errors.

Building upon these foundational advancements, RustAuditAI introduces a highly specialized, intra-procedural software quality intelligence framework designed to address the unique complexities of the Rust language ecosystem. While existing literature focuses heavily on binary vulnerability detection or syntax-focused parsing, this project shifts the paradigm toward a multi-dimensional quality assessment engine. By extending the classical Code Property Graph with a custom Function-Level Ownership Graph (FLOG), the framework maps Rust's precise variable lifetimes and move metrics into queryable graph spaces. By feeding these structural graph contexts into state-of-the-art reasoning layers, RustAuditAI unifies deterministic software engineering formulas (such as McCabe’s Cyclomatic Complexity) with context-aware, explainable AI recommendation streams. This integration delivers an actionable, mathematically grounded, and developer-friendly code intelligence tool tailored for modern secure systems development.
